\begin{ouverture}[probabilites]{19}{Hasard et probabilités}{Partie 4 · Données, hasard et programmation}
{Expérience · Issue · Évènement · Équiprobabilité · Fréquence}
\objectif{Décrire une expérience aléatoire et ses issues.}
\objectif{Attribuer une probabilité dans des cas équiprobables.}
\objectif{Comparer probabilité et fréquence observée.}
\prerequis{Fractions et fréquences.}
\end{ouverture}
\pagestyle{otMain}\essentiel\fontsize{11}{15}\selectfont
\phantomsection\label{cours-19-1}
\essSection{1}{Décrire ce qui peut arriver}
\begin{definition}Une expérience aléatoire est une expérience dont on ne connaît pas à l’avance le résultat. Une issue est un résultat possible. Un évènement réunit une ou plusieurs issues. Pour un dé, « obtenir un nombre pair » réunit 2, 4 et 6.\end{definition}
\label{fig-cours-19-1-cours-probabilites}\figureGeo{cours-probabilites}

\begin{propriete}Une probabilité est comprise entre 0 et 1. Un évènement impossible a une probabilité 0 ; un évènement certain une probabilité 1. Une chance sur quatre correspond à $1/4=0{,}25=25\,\%$.\end{propriete}
\begin{attention}Une probabilité très faible n’est pas zéro. Un résultat très probable n’est pas certain. Le hasard ne permet pas de prédire chaque résultat individuel.\end{attention}
\clearpage
\phantomsection\label{cours-19-2}
\essSection{2}{Compter et expérimenter}
\begin{definition}On parle d’équiprobabilité lorsque toutes les issues retenues ont la même probabilité. Pour un dé équilibré, les six faces sont équiprobables.\end{definition}
\begin{propriete}En cas d’équiprobabilité, la probabilité d’un évènement est le nombre d’issues favorables divisé par le nombre total d’issues. Pour un dé équilibré, obtenir plus de 4 a pour probabilité $2/6=1/3$.\end{propriete}
\begin{exemple}Une urne contient trois boules rouges et deux bleues, indiscernables au toucher. Chaque boule a la même chance d’être tirée, mais les deux couleurs n’ont pas la même chance : $3/5$ pour rouge et $2/5$ pour bleu.\end{exemple}
\begin{savoirfaire}[Confronter au réel]Répéter la même expérience en remettant les conditions de départ, noter chaque résultat puis calculer effectifs et fréquences. Une fréquence observée n’est pas nécessairement égale à la probabilité du modèle, surtout sur peu d’essais.\end{savoirfaire}

\clearpage\rubrique{Découverte}
\begin{blocLivre}{ex-19-1}\begin{exercice}[1]{Les issues d’un dé}
On lance un dé équilibré à six faces numérotées de 1 à 6. Lister les issues.
\end{exercice}\end{blocLivre}
\begin{blocLivre}{ex-19-2}\begin{exercice}[1]{Une pièce équilibrée}
Quelle est la probabilité d’obtenir pile en lançant une pièce équilibrée ? Donner trois écritures.
\end{exercice}\end{blocLivre}
\begin{blocLivre}{ex-19-3}\begin{exercice}[1]{Une échelle de probabilité}
Placer sur l’échelle de probabilité les évènements : A « obtenir 8 avec le dé précédent » ; B « obtenir un nombre inférieur à 7 avec ce dé » ; C « obtenir pile avec une pièce équilibrée » ; D « ne pas trouver la bonne combinaison au loto » ; E « obtenir dix fois de suite 1 avec le dé ».
\label{fig-ex-19-3}\figureGeo{proba-3}
\end{exercice}\end{blocLivre}
\begin{blocLivre}{ex-19-4}\begin{exercice}[1]{Une face donnée}
Quelle est la probabilité d’obtenir 5 sur un dé équilibré à six faces ?
\end{exercice}\end{blocLivre}
\begin{blocLivre}{ex-19-5}\begin{exercice}[1]{Lire une chance sur quatre}
Une chance sur quatre correspond-elle à 4 \%, 25 \% ou 40 \% ?
\end{exercice}\end{blocLivre}

\rubrique{Entraînement}
\begin{blocLivre}{ex-19-6}\begin{exercice}[2]{Un évènement à plusieurs issues}
Pour un dé équilibré, calculer la probabilité d’obtenir un nombre pair.
\end{exercice}\end{blocLivre}
\begin{blocLivre}{ex-19-7}\begin{exercice}[2]{Une urne}
Une urne contient quatre boules rouges, trois bleues et une verte, toutes indiscernables au toucher. Donner la probabilité de tirer bleu.
\label{fig-ex-19-7}\figureGeo{proba-7}
\end{exercice}\end{blocLivre}
\begin{blocLivre}{ex-19-8}\begin{exercice}[2]{Une roue régulière}
Une roue équilibrée est partagée en huit secteurs égaux, numérotés de 1 à 8. Calculer la probabilité d’un numéro supérieur à 5.
\label{fig-ex-19-8}\figureGeo{proba-8}
\end{exercice}\end{blocLivre}
\begin{blocLivre}{ex-19-9}\begin{exercice}[2]{Des jetons identifiés}
On tire un jeton au hasard parmi R1, R2, R3, B1, B2 et V1, indiscernables au toucher. Calculer la probabilité d’obtenir un jeton portant le numéro 1.
\label{fig-ex-19-9}\figureGeo{proba-9}
\end{exercice}\end{blocLivre}
\begin{blocLivre}{ex-19-10}\begin{exercice}[2]{Fréquence observée}
Une pièce donne 18 piles en 40 lancers. Calculer la fréquence de pile. Cela prouve-t-il que la pièce n’est pas équilibrée ?
\end{exercice}\end{blocLivre}

\rubrique{Problèmes}
\begin{blocLivre}{ex-19-11}\begin{exercice}[3]{Préparer une expérience}
Décrire une expérience pour observer la fréquence de rouge dans une urne avec deux rouges et trois bleues. Pourquoi faut-il remettre la boule ?
\label{fig-ex-19-11}\figureGeo{proba-11}
\end{exercice}\end{blocLivre}
\begin{blocLivre}{ex-19-12}\begin{exercice}[3]{Comparer deux jeux}
A fait gagner sur une face marquée parmi six faces équilibrées. B fait gagner sur trois jetons parmi dix jetons équiprobables. Quel jeu offre la plus grande probabilité de gagner ?
\end{exercice}\end{blocLivre}
\begin{blocLivre}{ex-19-13}\begin{exercice}[3]{Construire une urne}
Une urne de douze boules doit donner une probabilité de rouge égale à $1/3$. Combien de rouges faut-il ? Les autres boules seront bleues.
\label{fig-ex-19-13}\figureGeo{proba-13}
\end{exercice}\end{blocLivre}
\begin{blocLivre}{ex-19-14}\begin{exercice}[3]{Une fausse compensation}
Une pièce équilibrée donne pile cinq fois de suite. Lina affirme que le prochain lancer doit donner face pour compenser. Peut-elle le garantir ?
\end{exercice}\end{blocLivre}
\begin{blocLivre}{ex-19-15}\begin{exercice}[3]{Deux relevés différents}
On observe 6 succès sur 10 essais puis 46 succès sur 100 essais dans une expérience dont le modèle donne $1/2$. Comparer les fréquences et leurs écarts à $1/2$. Peut-on exiger que chaque série plus longue soit plus proche ?
\end{exercice}\end{blocLivre}

\clearpage
